\documentclass[11pt,a4paper]{article}
% Préambule commun aux documents de PythonIPM/documentation.
% Compilation : tectonic <fichier>.tex
% tectonic compile avec XeTeX, qui lit l'UTF-8 nativement : surtout PAS inputenc ni fontenc,
% faits pour pdfLaTeX. Ce sont eux qui décodaient mal les caractères multi-octets — « cœur »
% ressortait en « cur », et les tirets cadratins en caractères parasites.
\usepackage{lmodern}
\usepackage[french]{babel}
\usepackage{microtype}
\usepackage[a4paper,margin=2.4cm]{geometry}
\usepackage{amsmath,amssymb}
\usepackage{booktabs}
\usepackage{longtable}
\usepackage{array}
\usepackage{ragged2e}
\usepackage{xcolor}
\usepackage{tcolorbox}
\usepackage{enumitem}
\usepackage{fancyhdr}
\usepackage[hidelinks]{hyperref}

\definecolor{encadre}{HTML}{F4F6F8}
\definecolor{bordure}{HTML}{4A6FA5}
\definecolor{alerte}{HTML}{B03A2E}

% Encadré « à retenir » : la lecture non technique du passage qui précède
\newtcolorbox{retenir}[1][]{
  colback=encadre, colframe=bordure, boxrule=0.6pt, arc=2pt,
  left=8pt, right=8pt, top=6pt, bottom=6pt,
  fonttitle=\bfseries, title={#1}
}

% Encadré d'avertissement : un point ouvert, un écart assumé
\newtcolorbox{attention}[1][]{
  colback=white, colframe=alerte, boxrule=0.6pt, arc=2pt,
  left=8pt, right=8pt, top=6pt, bottom=6pt,
  % le bandeau de titre a déjà la couleur d'alerte en fond : un titre `coltitle=alerte`
  % écrirait du rouge sur du rouge. On garde le blanc par défaut, comme l'encadré « retenir ».
  fonttitle=\bfseries, title={#1}
}

\newcommand{\col}[1]{\texttt{#1}}
% \ensuremath et non \textsubscript : \ipm est employé aussi bien dans le texte que dans les
% formules, et « M\textsubscript{0} » placé en mode mathématique est une erreur — silencieuse en
% PDF, mais qui fait échouer la conversion vers Word.
\newcommand{\ipm}{\ensuremath{M_{0}}}

\setlist[itemize]{itemsep=2pt, topsep=4pt}
\setlist[description]{style=nextline, leftmargin=1.2cm, itemsep=5pt}

% La source du document apparaît en tête de chaque page. Les documents RGPH la redéfinissent
% par \renewcommand{\source}{...} AVANT \begin{document}.
\newcommand{\source}{EHCVM 2021}

\pagestyle{fancy}
\fancyhf{}
\fancyhead[L]{\small\itshape IPM Côte d'Ivoire --- \source}
\fancyhead[R]{\small\thepage}
\renewcommand{\headrulewidth}{0.3pt}

\renewcommand{\arraystretch}{1.15}

\renewcommand{\source}{RGPH 2021}

\title{\bfseries L'indice de pauvreté multidimensionnelle de la Côte d'Ivoire\\[4pt]
\large Méthodologie de calcul sur le RGPH 2021}
\author{Pipeline \col{PythonIPM}}
\date{Septembre 2026}

\begin{document}
\maketitle

\begin{abstract}
\noindent Ce document décrit la construction de l'indice de pauvreté multidimensionnelle (IPM) de
la Côte d'Ivoire à partir du Recensement général de la population et de l'habitation (RGPH) 2021.
Il est le pendant du document consacré à l'EHCVM 2021 : même méthode d'Alkire-Foster, mêmes
conventions de calcul, même chaîne de traitement. Ce qui change tient aux données — un recensement
exhaustif de 28 millions d'individus, au questionnaire beaucoup plus court qu'une enquête ménages.
Le document se concentre donc sur ce qui est propre au recensement : les indicateurs qu'il permet
de mesurer et ceux qu'il interdit, les approximations retenues, la réconciliation avec la
production Stata antérieure, et la variante harmonisée qui rend les deux sources comparables.
\end{abstract}

\begin{retenir}[Les résultats en six lignes]
\begin{itemize}
  \item \textbf{38,8~\%} de la population ivoirienne vit dans un ménage multidimensionnellement
        pauvre, soit 11,4 millions de personnes --- IC 95~\% : [38,4~\% ; 39,3~\%].
  \item Ces personnes cumulent en moyenne \textbf{42,9~\%} des privations pondérées possibles.
  \item L'indice qui combine les deux vaut $\ipm = \mathbf{0{,}1664}$, IC 95~\%
        [0,1645 ; 0,1684].
  \item Le milieu rural rassemble \textbf{84,7~\%} de la pauvreté pour 45,8~\% de la population ;
        l'urbain en porte 5,7~\% pour 44,8~\% de la population.
  \item L'écart entre régions va de 0,007 (Abidjan) à 0,365 (Bounkani), soit un rapport de
        un à cinquante-trois.
  \item Sur le socle commun avec l'EHCVM, les deux sources se rejoignent :
        $\ipm = 0{,}2680$ au recensement contre $0{,}2737$ à l'enquête, soit 2~\% d'écart
        (section~\ref{sec:harmonise}).
\end{itemize}
\end{retenir}

\tableofcontents
\bigskip

\section{Pourquoi calculer l'IPM sur un recensement}

L'EHCVM interroge 12\,965 ménages. C'est assez pour un indice national et pour les 33 régions,
mais la précision se dégrade dès qu'on descend plus bas : beaucoup de sous-préfectures n'y sont
représentées que par une poignée de grappes.

Le RGPH, lui, dénombre \textbf{tous} les ménages du pays. Il ouvre donc ce que l'enquête ne peut
pas offrir : une carte de la pauvreté multidimensionnelle à la maille de la sous-préfecture et de
la commune, avec des effectifs qui la rendent publiable presque partout.

La contrepartie est le questionnaire. Un recensement doit être administré à trente millions de
personnes : il est court. Il ne pose aucune question de santé, aucune sur le nombre de pièces du
logement, aucune sur le nombre d'années d'études effectivement accomplies. Trois indicateurs de
l'IPM national y sont donc impossibles, un quatrième aussi, et deux autres ne peuvent être
qu'approchés.

\begin{retenir}[Le partage des rôles]
L'EHCVM donne l'indice \emph{le plus complet} --- seize indicateurs, dont la santé et
l'alimentation. Le RGPH donne l'indice \emph{le plus fin géographiquement} --- 519
sous-préfectures et communes contre 442 sur un échantillon. Les deux ne se remplacent pas, ils se
complètent.
\end{retenir}

\section{Les données}

Le RGPH 2021 couvre \textbf{5\,616\,487 ménages} et \textbf{27\,984\,405 individus} recensés, soit
\textbf{29\,276\,660 personnes} une fois la correction de couverture appliquée. L'extrait utilisé
tient en trois fichiers.

\begin{center}
\begin{tabular}{@{}llr>{\RaggedRight}p{5.6cm}@{}}
\toprule
\textbf{Fichier} & \textbf{Une ligne =} & \textbf{Lignes} & \textbf{Ce qu'on y prend} \\
\midrule
\col{IPM\_Data\_191125} & un individu & 27\,984\,405 & âge, éducation, emploi, état civil \\
\col{IPM\_Data\_110924\_men} & un ménage (le chef) & 5\,528\,535 & logement, eau, énergie,
  sanitaires, biens \\
\col{MORTALITE\_RP2021} & un décès & 269\,901 & dimension Santé \\
\bottomrule
\end{tabular}
\end{center}

\subsection{Pourquoi deux fichiers plutôt qu'un}

Le fichier individus porte lui aussi les variables de logement et de biens d'équipement. Mais
elles n'y sont renseignées que pour \textbf{un tiers des ménages} : une fusion incomplète en amont
les a perdues pour les deux autres tiers. Les lire là reviendrait à priver de conditions de vie
deux ménages sur trois.

Elles sont donc lues dans le fichier ménages, où elles sont complètes, et rattachées au ménage par
l'\col{INDIV\_ID} de son chef. Restent \textbf{87\,952 ménages}, soit 1,6~\%, présents dans le
fichier individus et absents du fichier ménages : conformément à la convention du pipeline, ils ne
sont privés sur aucun indicateur de conditions de vie.

\subsection{Les trois clés de jointure}

\begin{description}
  \item[\col{ID\_Menage}] rassemble les individus d'un même ménage, à l'intérieur du fichier
    individus.
  \item[\col{INDIV\_ID} du chef] rattache le ménage à sa ligne du fichier ménages.
  \item[la clé géographique complète] (\col{REGION} \dots{} \col{P10}, douze variables) rattache le
    ménage à ses décès. C'est le seul lien possible : le fichier de mortalité ne porte ni
    \col{INDIV\_ID} ni \col{ID\_Menage}.
\end{description}

\subsection{La pondération}

Un recensement se veut exhaustif, mais aucun ne l'est parfaitement. L'Enquête post-censitaire
mesure le taux de couverture atteint ; la variable \col{EW} en porte la correction, comprise entre
1,020 et 1,057 selon la strate.

Appliquée, elle donne une population de \textbf{29,28 millions}, à 0,4~\% du chiffre publié du RGPH
2021 (29\,389\,150). C'est le contrôle qui valide le choix de la variable.

\begin{attention}[Ne pas confondre \col{EW} et \col{PROJ24}]
La base porte une seconde pondération, \col{PROJ24}, dite « de ratissage », comprise entre 1,109 et
1,153. Elle projette les effectifs à 2024 et donnerait 31,9 millions de personnes. Elle n'est pas
retenue : l'IPM publié ici décrit la Côte d'Ivoire \emph{de 2021}.
\end{attention}

\subsection{La grappe : une zone de dénombrement}

Les intervalles de confiance de la méthode agrègent la variance à l'unité primaire de sondage.
Un recensement n'en a pas, à proprement parler : il n'y a pas de tirage. On retient la
\textbf{zone de dénombrement} --- \col{SOUSPREFID} $\times$ \col{P\_ZC} $\times$ \col{P05} --- soit
\textbf{28\,329 unités}, qui est l'unité de collecte et la maille à laquelle les erreurs de
dénombrement se corrèlent.

\begin{retenir}[Comment lire un intervalle de confiance sur un recensement]
Il ne mesure pas une incertitude d'échantillonnage --- il n'y en a pas --- mais la dispersion du
phénomène entre zones de dénombrement. Il reste utile pour juger si une sous-préfecture au chiffre
extrême est réellement extrême ou portée par quelques zones atypiques. Il est, par construction,
bien plus étroit que celui de l'EHCVM.
\end{retenir}

\subsection{Unité de construction, unité de comptage}

Inchangé par rapport à l'EHCVM : le score se calcule par ménage, l'indice se lit en personnes.
Chaque ménage $i$ compte pour
\[
n_i \;=\; \col{taille\_menage}_i \times \col{ponderation\_menage}_i .
\]

\section{Les treize indicateurs}
\label{sec:indicateurs}

Quatre dimensions équipondérées à 0,25, comme dans l'IPM national. Le poids d'une dimension se
partage à parts égales entre ses indicateurs : une dimension à un seul indicateur lui donne donc
tout son poids.

\begin{center}
\small
\begin{tabular}{@{}l>{\RaggedRight}p{5.0cm}r r@{}}
\toprule
\textbf{Dimension} & \textbf{Indicateur} & \textbf{Poids $w_j$} & \textbf{Privés} \\
\midrule
Éducation (0,25) & Fréquentation scolaire & 0,0625 & 35,6~\% \\
                 & Année de scolarité & 0,0625 & 63,2~\% \\
                 & Alphabétisation & 0,0625 & 70,4~\% \\
                 & Déclaration d'état civil & 0,0625 & 15,2~\% \\
\addlinespace
Santé (0,25) & Mortalité dans le ménage & 0,2500 & 1,1~\% \\
\addlinespace
Emploi (0,25) & Chômage & 0,1250 & 5,9~\% \\
              & Emploi agricole de subsistance & 0,1250 & 46,0~\% \\
\addlinespace
Conditions de vie (0,25) & Électricité & 0,0417 & 11,5~\% \\
                         & Logement & 0,0417 & 25,0~\% \\
                         & Eau potable & 0,0417 & 14,3~\% \\
                         & Énergie de cuisson & 0,0417 & 62,7~\% \\
                         & Toilette & 0,0417 & 45,1~\% \\
                         & Biens d'équipement & 0,0417 & 18,2~\% \\
\bottomrule
\end{tabular}
\end{center}

\noindent « Privés » est la part de la \emph{population} vivant dans un ménage privé sur
l'indicateur.

\subsection{Ce que le recensement mesure comme l'enquête}

Sept indicateurs sont transposés sans perte : fréquentation scolaire, alphabétisation, état civil,
chômage, électricité, énergie de cuisson et logement. Seuls les codes de nomenclature changent.

Un indicateur y gagne même en fidélité. La définition PNUD des biens d'équipement compte
« radio, télévision, téléphone, ordinateur, charrette, vélo, moto ou réfrigérateur ». La section 12
de l'EHCVM ne propose pas la charrette : l'indice y porte sur sept biens. Le recensement, lui,
demande le « véhicule à traction animale » : les \textbf{huit biens y sont tous}.

\subsection{Les deux approximations}

\begin{description}
  \item[Année de scolarité.] L'énoncé national dit « aucun membre de 17-95 ans n'a complété dix
    années d'études ». Le recensement ne demande ni les années d'études ni la classe atteinte : il
    ne demande que le \textbf{diplôme le plus élevé}. Le seuil est donc porté sur le diplôme --- le
    BEPC ou plus vaut dix années accomplies. L'approximation est bonne à l'échelle du pays : 63,2~\%
    de privés au recensement contre 62,9~\% à l'enquête, où les années d'études sont reconstruites.
    Elle l'est moins individuellement : on peut avoir fait la troisième sans obtenir le BEPC.

  \item[Emploi agricole de subsistance.] L'énoncé national vise le chef de ménage dont l'activité se
    limite à l'agriculture sur son propre champ. Le recensement ne demande pas « sur son propre
    champ ». On retient donc le chef occupé dans une \textbf{branche agricole} (divisions 01 à 03 de
    la NAEMA) \textbf{à son compte ou comme aide familial} --- ce qui exclut le salariat,
    l'apprentissage et les employeurs. 46,0~\% de privés, contre 40,5~\% à l'enquête.
\end{description}

\subsection{Les deux indicateurs amputés d'un critère}

L'eau potable de la définition PNUD est privée si la source n'est pas améliorée au sens des ODD
\emph{ou} si elle est à trente minutes ou plus aller-retour. La toilette est privée si les
installations ne sont pas améliorées \emph{ou} si elles sont partagées. Le recensement ne demande
ni le temps de trajet, ni le partage : seul le premier critère s'applique dans les deux cas.

L'effet est massif sur la toilette --- 45,1~\% de privés au recensement contre 67,2~\% à l'enquête,
où le partage des sanitaires pèse lourd en milieu urbain. Il est plus modéré sur l'eau : 14,3~\%
contre 20,9~\%.

\subsection{Les quatre indicateurs impossibles}

Le recensement ne pose \textbf{aucune question de santé} : ni assurance maladie, ni module de
sécurité alimentaire, ni recours aux soins. Il ne demande pas non plus le \textbf{nombre de pièces}
du logement, ce qui interdit la promiscuité. Quatre indicateurs de l'IPM national disparaissent
donc : assurance maladie, insécurité alimentaire, renoncement aux soins et promiscuité.

\subsection{L'indicateur propre au recensement : la mortalité}
\label{sec:mortalite}

La dimension Santé serait vide sans le fichier de mortalité, qui recense les décès survenus dans le
ménage au cours des douze mois précédant le recensement. On retient le décès d'un
\textbf{enfant de moins de 18 ans} : c'est la borne de l'IPM mondial, et celle du traitement Stata
antérieur du RGPH 2021.

\begin{attention}[Une privation rare, qui porte pourtant le quart du poids]
Seuls \textbf{1,1~\%} des ménages sont privés sur cet indicateur, et il est seul dans sa dimension :
son poids $w_j$ vaut 0,25. La dimension Santé ne contribue donc qu'à \textbf{1,6~\%} de \ipm{}
pour un poids nominal de 25~\%, et le score d'un ménage ne peut pratiquement pas dépasser 0,75.
C'est la première raison pour laquelle l'IPM RGPH complet est structurellement plus bas que l'IPM
EHCVM, et pourquoi les deux ne se comparent pas en niveau.

Deux facteurs se cumulent. L'IPM mondial retient une fenêtre de \textbf{cinq ans} là où le
recensement n'en couvre qu'\textbf{une}. Et l'âge au décès est inconnu pour 27,1~\% des décès
enregistrés, qui ne peuvent donc pas être comptés.
\end{attention}

\section{Le calcul}

Identique à celui de l'EHCVM --- le code des étapes 02 à 05 est littéralement le même fichier. On
en rappelle ici les quatre temps ; le détail des formules figure dans le document méthodologique
de l'EHCVM.

\begin{enumerate}
  \item \textbf{Matrice de privation $g^0$.} Le vecteur de seuils $z$ est appliqué aux situations
        brutes : une colonne 0/1 par indicateur, 1 = privé.
  \item \textbf{Pondération et score.} $c_i = \sum_j w_j\, g^0_{ij}$, entre 0 et 1.
  \item \textbf{Censure au seuil $k = 1/3$.} Les privations des ménages non pauvres sont remises à
        zéro, ce qui rend \ipm{} décomposable par indicateur.
  \item \textbf{Indices.}
        \[
        H = \frac{\sum_i n_i \mathbf{1}(c_i \geq k)}{\sum_i n_i}, \qquad
        A = \frac{\sum_i n_i c_i(k)}{\sum_i n_i \mathbf{1}(c_i \geq k)}, \qquad
        \ipm = H \times A .
        \]
\end{enumerate}

Trois variantes sont produites à chaque exécution : l'IPM national à quatre dimensions, l'IPM
international à trois dimensions (sans l'Emploi), et l'IPM harmonisé décrit en
section~\ref{sec:harmonise}.

\section{Résultats}

\subsection{L'indice national}

\begin{center}
\begin{tabular}{@{}lrrr@{}}
\toprule
\textbf{Variante} & $H$ & $A$ & \ipm \\
\midrule
IPM national --- 4 dimensions, 13 indicateurs & 38,8~\% & 0,4287 & \textbf{0,1664} \\
IPM international --- 3 dimensions, sans l'Emploi & 36,1~\% & 0,4301 & 0,1553 \\
IPM harmonisé --- socle commun EHCVM/RGPH & 52,0~\% & 0,5157 & 0,2680 \\
\bottomrule
\end{tabular}
\end{center}

\noindent Sur l'IPM national, \textbf{20,7~\%} de la population est vulnérable ($0{,}2 \le c_i <
1/3$) et \textbf{7,8~\%} en pauvreté sévère ($c_i \ge 0{,}5$).

\subsection{Le milieu de résidence}

Le recensement distingue trois milieux là où l'enquête n'en code que deux.

\begin{center}
\begin{tabular}{@{}lrrrrr@{}}
\toprule
\textbf{Milieu} & \textbf{Part pop.} & $H$ & $A$ & \ipm & \textbf{Contribution} \\
\midrule
Urbain & 44,8~\% & 5,4~\% & 0,3943 & 0,0212 & 5,7~\% \\
Semi-urbain & 9,4~\% & 41,0~\% & 0,4157 & 0,1705 & 9,7~\% \\
Rural & 45,8~\% & 71,1~\% & 0,4328 & 0,3075 & \textbf{84,7~\%} \\
\bottomrule
\end{tabular}
\end{center}

C'est le résultat le plus net du calcul. Le rural porte \textbf{84,7~\%} de la pauvreté
multidimensionnelle pour 45,8~\% de la population ; l'urbain en porte 5,7~\% pour une part de
population presque identique. Le rapport des \ipm{} entre rural et urbain est de \textbf{un à
quinze} (0,3075 contre 0,0212), contre un à trois dans l'EHCVM (0,4354 contre 0,1525) : le
recensement fait ressortir un contraste que l'échantillon lissait.

Le semi-urbain, que l'enquête ne distingue pas, occupe une position intermédiaire cohérente --- et
c'est le milieu où la \textbf{vulnérabilité} est la plus forte (34,2~\% contre 17,7~\% en urbain) :
une population nombreuse y est juste sous le seuil.

\subsection{Le territoire}

\begin{center}
\small
\begin{tabular}{@{}lrrrr@{}}
\toprule
\textbf{Région} & \textbf{Part pop.} & $H$ & \ipm & \textbf{Contribution} \\
\midrule
\multicolumn{5}{@{}l}{\itshape Les cinq plus faibles} \\
District autonome d'Abidjan & 21,5~\% & 1,8~\% & 0,0069 & 0,9~\% \\
District autonome de Yamoussoukro & 1,4~\% & 18,7~\% & 0,0755 & 0,6~\% \\
Sud-Comoé & 2,7~\% & 28,1~\% & 0,1181 & 1,9~\% \\
Gbêkê & 4,6~\% & 30,2~\% & 0,1279 & 3,5~\% \\
La Mé & 2,5~\% & 32,2~\% & 0,1324 & 2,0~\% \\
\addlinespace
\multicolumn{5}{@{}l}{\itshape Les cinq plus élevées} \\
Gbôklé & 1,6~\% & 65,1~\% & 0,2804 & 2,6~\% \\
Béré & 1,7~\% & 66,9~\% & 0,2843 & 2,9~\% \\
Bafing & 0,9~\% & 69,5~\% & 0,3139 & 1,7~\% \\
Folon & 0,5~\% & 76,2~\% & 0,3275 & 1,0~\% \\
Bounkani & 1,5~\% & 76,5~\% & 0,3654 & 3,2~\% \\
\bottomrule
\end{tabular}
\end{center}

L'écart va de 0,0069 à 0,3654, soit un rapport de \textbf{un à cinquante-trois}. Abidjan tire seule
la moyenne nationale vers le bas : 21,5~\% de la population du pays, moins de 1~\% de sa pauvreté
multidimensionnelle.

À la maille du département, l'écart va d'Abidjan (\ipm{} = 0,0069) à Téhini (0,4191), et à celle de
la sous-préfecture, du Plateau (0,0015) à Santa de Ouaninou (0,4752), où 92,6~\% de la population
est multidimensionnellement pauvre. C'est cette granularité qui justifie le calcul sur recensement.

\subsection{La taille du ménage}

\begin{center}
\begin{tabular}{@{}lrrrr@{}}
\toprule
\textbf{Taille} & \textbf{Part pop.} & $H$ & \ipm & \textbf{Contribution} \\
\midrule
1-2 personnes & 8,4~\% & 20,9~\% & 0,0825 & 4,2~\% \\
3-4 & 17,8~\% & 28,8~\% & 0,1187 & 12,7~\% \\
5-6 & 22,3~\% & 34,9~\% & 0,1491 & 19,9~\% \\
7-9 & 22,2~\% & 40,5~\% & 0,1747 & 23,3~\% \\
10 et plus & 29,3~\% & 51,8~\% & 0,2266 & \textbf{39,9~\%} \\
\bottomrule
\end{tabular}
\end{center}

La relation est monotone et forte. Elle doit se lire avec précaution : un grand ménage a
mécaniquement plus de chances de compter au moins un enfant non scolarisé ou au moins un chômeur,
puisque plusieurs indicateurs sont formulés en « au moins un membre ». Une partie de la pente est
donc un artefact de la définition, pas un fait sur les grands ménages.

\subsection{Ce que pèse réellement chaque dimension}

\begin{center}
\begin{tabular}{@{}lrr@{}}
\toprule
\textbf{Dimension} & \textbf{Poids nominal} & \textbf{Contribution à \ipm} \\
\midrule
Éducation & 25~\% & 40,6~\% \\
Conditions de vie & 25~\% & 29,2~\% \\
Emploi & 25~\% & 28,6~\% \\
Santé & 25~\% & \textbf{1,6~\%} \\
\bottomrule
\end{tabular}
\end{center}

Les quatre indicateurs qui portent l'indice sont l'emploi agricole de subsistance (26,8~\%),
l'alphabétisation (13,6~\%), l'année de scolarité (13,4~\%) et l'énergie de cuisson (9,5~\%). Le
chômage (1,8~\%) et la mortalité (1,6~\%) n'y pèsent presque rien, alors qu'ils portent ensemble
0,375 de poids nominal.

\begin{attention}[L'emploi de subsistance domine l'indice]
Cet indicateur est à la fois le \textbf{premier contributeur} à \ipm{} (26,8~\%) et celui qui repose
sur l'\textbf{approximation la plus lâche} (section~\ref{sec:indicateurs}). Deux effets se
conjuguent : la dimension Emploi ne compte que deux indicateurs, donc chacun pèse 0,125 ; et le
chômage ne prive presque personne, laissant l'emploi de subsistance porter seul la dimension.

C'est le point à réexaminer en priorité. Si la substitution de « branche agricole $\times$
indépendant ou aide familial » à « son propre champ » est jugée trop large, c'est plus du quart de
l'indice qui se déplace.
\end{attention}

\section{Réconciliation avec la production Stata antérieure}
\label{sec:stata}

Le RGPH 2021 avait déjà été traité en Stata (\col{rp21\_calcul\_des\_privations.do}), avec une
sortie conservée. Chaque règle de ce programme a été rejouée sur les données afin de séparer ce qui
relève d'un écart de définition de ce qui relèverait d'une erreur.

\begin{retenir}[Le résultat du contrôle]
\textbf{Aucune erreur.} Appliquée telle quelle, chaque règle Stata est reproduite à 0,1 point près
par le pipeline. Tous les écarts de résultat sont des choix de définition, documentés ci-dessous.
\end{retenir}

\begin{center}
\small
\begin{tabular}{@{}lrrr>{\RaggedRight}p{4.6cm}@{}}
\toprule
\textbf{Indicateur} & \textbf{Stata} & \textbf{Règle rejouée} & \textbf{Retenu} & \textbf{Écart dû à} \\
\midrule
Fréquentation scolaire & 30,10~\% & 30,14~\% & 35,60~\% & version en 6-14 ans ; le programme
  \emph{courant} dit 6-16, comme ici \\
Année de scolarité & 63,05~\% & 63,15~\% & 63,17~\% & --- \\
Alphabétisation & 30,10~\% & 33,15~\% & 70,43~\% & lecture de l'énoncé \\
Chômage & 2,58~\% & 5,27~\% & 5,92~\% & tranche 16-35 contre 17-40 \\
État civil & 9,65~\% & 9,73~\% & 15,22~\% & déclaration contre possession \\
Électricité & 11,59~\% & --- & 11,5~\% & --- \\
Eau potable & 14,43~\% & --- & 14,3~\% & --- \\
Énergie de cuisson & 63,40~\% & --- & 62,7~\% & --- \\
Logement & 25,25~\% & --- & 25,0~\% & --- \\
Biens d'équipement & 18,72~\% & --- & 18,2~\% & --- \\
Mortalité & 1,121~\% & --- & 1,1~\% & --- \\
Toilette & 33,09~\% & 46,55~\% & 45,1~\% & \textbf{non réconcilié} \\
\bottomrule
\end{tabular}
\end{center}

\subsection{Trois corrections apportées après cette comparaison}

Toutes dans le sens du programme Stata officiel.

\begin{enumerate}
  \item \textbf{Mortalité : moins de 18 ans}, et non moins de 5. C'est la borne du programme
        \emph{et} celle de l'IPM mondial. Le taux passe de 0,7~\% à 1,1~\%, soit exactement le
        1,121~\% de la sortie Stata.
  \item \textbf{Logement : matériaux définis par exclusion}, comme le programme. Est rudimentaire
        tout ce qui n'est ni ciment, carreau ou moquette (sol), ni tôle, béton ou tuile (toit), ni
        semi-dur, géobéton ou dur (mur). Cela ajoute le sol en bois et le mur en tôle, conformément
        à la classification DHS. Le taux passe de 24,0~\% à 25,0~\%, contre 25,25~\% en Stata.
  \item \textbf{Diplôme :} le code 22 « Autres à préciser » compte au-dessus du seuil, comme le
        programme (0,02 point).
\end{enumerate}

\subsection{Quatre écarts maintenus}

Le cadrage retenu est celui de l'EHCVM : là où les deux se contredisent, c'est l'énoncé national
qui tranche.

\begin{description}
  \item[Alphabétisation.] Le programme RGPH code « \textbf{aucun} membre alphabétisé » (33,2~\%) ;
    l'énoncé national dit « \textbf{un} membre ne sait pas lire ou écrire » (70,4~\%). L'énoncé
    littéral est appliqué, aux deux sources, par un drapeau unique dans le code : elles ne peuvent
    pas diverger sur ce point sans qu'on le décide.

  \item[État civil.] La question de l'EHCVM (1.05) est « dispose d'un acte de naissance ? ». La
    modalité RGPH « déclaré \emph{sans} acte de naissance » est donc une privation. Le programme
    Stata mesure la déclaration et non la possession, et ne compte privés que « non déclaré » et
    « ne sait pas ». La lecture retenue (15,2~\%) est aussi la plus proche de l'EHCVM (26,3~\%).

  \item[Chômage sur les 17-40 ans], énoncé du tableau national, contre 16-35 dans le programme
    RGPH. La \emph{variable} de statut, elle, a été validée : \col{Statut\_OQPtbb} ne diffère de la
    construction du programme --- les trois critères du BIT, sans emploi, en recherche, disponible
    --- que sur 18\,910 individus sur 28 millions, soit 0,07~\%. L'autre variable de la base,
    \col{Statut\_OQPtb}, donnerait 21,6~\% de ménages privés : elle n'intègre pas le reclassement
    des personnes déclarées sans activité mais en réalité occupées, et n'est pas retenue.

  \item[Emploi agricole de subsistance], absent de l'IPM RGPH officiel mais présent dans les seize
    indicateurs de l'EHCVM : il est donc calculé, par approximation.
\end{description}

\subsection{Un écart non réconcilié : la toilette}

La règle est identique de part et d'autre, le fichier source est le même, et pourtant la sortie
Stata donne 33,09~\% quand le pipeline donne 46,55~\% avant convention sur les manquants.

Aucune version du programme présente dans le dépôt ne produit 33~\%. La liste qui ajouterait la
chasse d'eau à l'air libre et vers un lieu inconnu donnerait 34,19~\%, soit l'ordre de grandeur de
la sortie Stata : celle-ci provient vraisemblablement d'un millésime antérieur du code ou des
données.

\begin{attention}[Point à confirmer auprès de l'INS]
La valeur retenue ici est celle du programme \emph{courant}. Elle est cohérente avec l'EHCVM, qui
donne 67,2~\% en ajoutant le critère de partage des sanitaires, absent du recensement. Mais
l'origine du 33,09~\% n'est pas établie et mérite d'être élucidée avant publication.
\end{attention}

\subsection{Ce qui n'est pas comparable, et pourquoi}

Le programme Stata construit un IPM à \textbf{cinq dimensions équipondérées à 0,2} --- Éducation,
Emploi, Conditions de vie, Santé et \textbf{Identification}, l'état civil y étant une dimension à
part entière --- et remplace les valeurs manquantes par une privation.

Le pipeline suit le cadrage EHCVM : \textbf{quatre dimensions à 0,25}, l'état civil rangé dans
l'Éducation, et un ménage dont aucune situation n'est renseignée compté non privé. Les
$H = 0{,}100$ et $\ipm = 0{,}045$ de la sortie Stata ne se comparent donc pas aux
$H = 0{,}388$ et $\ipm = 0{,}166$ produits ici.

\section{La variante harmonisée : comparer les deux sources}
\label{sec:harmonise}

Les IPM complets des deux sources ne se comparent pas en niveau : chacun exploite ce que sa source
mesure, et la dimension Santé du recensement --- la seule mortalité, privation à 1,1~\% --- absorbe
un quart du poids pour presque rien.

Le pipeline produit donc, pour chaque source, une variante restreinte au \textbf{socle commun} :
les douze indicateurs que l'EHCVM \emph{et} le RGPH savent mesurer. La dimension Santé n'en a aucun
--- l'enquête et le recensement n'ont pas un seul indicateur de santé en commun --- la variante
compte donc \textbf{trois dimensions à 1/3} : Éducation (4 indicateurs), Emploi (2) et Conditions
de vie (6).

\begin{center}
\begin{tabular}{@{}lrrr@{}}
\toprule
& $H$ & $A$ & \ipm \\
\midrule
EHCVM 2021 & 53,5~\% & 0,5113 & \textbf{0,2737} \\
RGPH 2021 & 52,0~\% & 0,5157 & \textbf{0,2680} \\
\bottomrule
\end{tabular}
\end{center}

\begin{retenir}[Le meilleur contrôle de cohérence disponible]
À indicateurs identiques, une enquête par sondage de 12\,965 ménages et un recensement exhaustif de
5,6 millions de ménages, collectés la même année par deux dispositifs différents, donnent
\textbf{0,2737 et 0,2680} : \textbf{2~\% d'écart}. L'incidence diffère de 1,5 point, l'intensité de
0,004.
\end{retenir}

La décomposition confirme l'accord indicateur par indicateur. Les contributions à \ipm{} se
suivent de près : emploi de subsistance 23,3~\% contre 26,6~\%, année de scolarité 14,2~\% contre
14,1~\%, alphabétisation 13,6~\% contre 14,5~\%, énergie de cuisson 10,1~\% contre 10,0~\%.

Les deux seuls écarts notables sont ceux qu'on attend : la \textbf{toilette} (9,6~\% contre 7,5~\%),
puisque le recensement ne mesure pas le partage des sanitaires, et l'\textbf{état civil} (7,2~\%
contre 4,5~\%).

\begin{attention}[Ce que la variante harmonisée ne dit pas]
Elle valide la \emph{cohérence} des deux sources, pas l'IPM national. Celui-ci reste l'indice à
seize indicateurs de l'EHCVM ($\ipm = 0{,}2863$). La variante harmonisée sert au contrôle et à la
comparaison ; elle n'a pas vocation à être publiée comme l'indice du pays.
\end{attention}

\section{Le traitement des valeurs manquantes}

La convention du pipeline est inchangée : \textbf{un ménage non concerné, ou dont la situation n'est
pas renseignée, n'est pas privé}.

Elle prend une importance particulière sur le recensement, où 1,6~\% des ménages n'ont pas de ligne
dans le fichier des conditions de vie. Une règle formulée par exclusion --- « privé si la source
d'éclairage n'est pas parmi les modalités adéquates » --- rendrait mécaniquement privés tous ces
ménages, puisqu'une valeur manquante n'appartient à aucune liste.

Le pipeline pose donc un garde-fou explicite : si \textbf{aucune} des situations qui composent un
indicateur n'est renseignée, le ménage n'est pas privé sur cet indicateur. Il reste privé s'il est
seulement partiellement renseigné --- un ménage sans toilettes n'a pas de réponse à la question du
partage, mais son type de sanitaire suffit à le classer.

\begin{attention}[Une divergence de convention avec le programme Stata]
Le programme Stata fait l'inverse : il remplace tout manquant par une privation. Sur les 1,6~\% de
ménages concernés, les deux conventions donnent des résultats opposés. Celle du pipeline traite le
défaut d'appariement comme ce qu'il est --- une donnée absente --- et non comme une privation
constatée.
\end{attention}

\section{Le pipeline}

Le calcul est mené par la même chaîne que l'EHCVM. La source est portée par la variable
d'environnement \col{IPM\_SOURCE}, qui détermine le dossier de données, la clé du ménage, les
variables de désagrégation disponibles et le suffixe de tous les fichiers produits.

\begin{center}
\small
\begin{tabular}{@{}l>{\RaggedRight}p{5.4cm}>{\RaggedRight}p{5.0cm}@{}}
\toprule
\textbf{Étape} & \textbf{Fichier} & \textbf{Sortie} \\
\midrule
01 & \col{01\_rgph\_preconstruction\_\dots} & situations brutes, 5\,616\,487 $\times$ 29 \\
02 & \col{02\_construction\_\dots} & 13 indicateurs 0/1 \\
03 & \col{03\_matrice\_privations\_\dots} & $w_j$, $g^0$ pondérée, score $c_i$ \\
04 & \col{04\_matrice\_privations\_\dots} & censure au seuil $k$ \\
05 & \col{05\_indices\_ipm} & $H$, $A$, \ipm{}, désagrégations, contributions \\
\bottomrule
\end{tabular}
\end{center}

Seule l'étape 01 est propre au recensement : elle lit des bases sans rapport avec celles de
l'enquête. À partir de l'étape 02, le code est strictement commun --- il ne travaille plus que sur
des colonnes de situation normalisées, et ne sait pas de quelle source elles viennent. C'est ce qui
garantit que les deux IPM sont calculés par la même méthode, et non par deux implémentations qui
se ressemblent.

\begin{verbatim}
python pipeline/orchestrateur.py                  # les deux IPM
python pipeline/orchestrateur.py --source rgph    # le RGPH seul
python pipeline/orchestrateur.py --check          # les auto-contrôles
\end{verbatim}

L'étape 01 lit le fichier individus par morceaux de trois millions de lignes : vingt-huit millions
de lignes ne tiennent pas en mémoire. Un ménage coupé entre deux morceaux ne pose pas de problème,
chaque morceau étant agrégé séparément avant que les agrégats ne soient resommés par ménage. La
chaîne complète s'exécute en un peu plus de deux minutes et écrit 2,2~Go de tables de travail.

Chaque étape dispose d'un auto-contrôle qui rejoue sa logique sur un mini-jeu de données aux
résultats connus, et chaque exécution vérifie par assertion que $\ipm = H \times A$, que les
contributions somment à 1, et que les \ipm{} de groupe pondérés par leur part de population
redonnent l'\ipm{} national.

\section{Limites}

\begin{enumerate}
  \item \textbf{La dimension Santé ne tient qu'à un fil.} Un indicateur, une privation à 1,1~\%,
        une fenêtre de douze mois au lieu de cinq ans, et 27,1~\% d'âges au décès inconnus. Elle
        porte 25~\% du poids nominal et 1,6~\% de l'indice
        (section~\ref{sec:mortalite}).
  \item \textbf{L'emploi de subsistance domine l'indice} tout en reposant sur l'approximation la
        plus lâche : 26,8~\% de \ipm{}.
  \item \textbf{L'année de scolarité est mesurée par le diplôme}, ce qui est juste en moyenne mais
        faux individuellement.
  \item \textbf{L'eau et la toilette perdent leur second critère}, ce qui abaisse mécaniquement les
        deux taux par rapport à l'enquête.
  \item \textbf{Le taux de privation de la toilette n'est pas réconcilié} avec la production
        antérieure (section~\ref{sec:stata}).
  \item \textbf{Le sexe du chef de ménage n'est pas dans l'extrait}, ce qui prive d'une
        désagrégation disponible sur l'EHCVM.
  \item \textbf{Les intervalles de confiance ne mesurent pas une incertitude d'échantillonnage}
        mais une dispersion entre zones de dénombrement. Dix-neuf sous-préfectures sur 519
        dépassent un coefficient de variation de 20~\% et ne devraient pas être publiées telles
        quelles.
\end{enumerate}

\end{document}
